\section{Schwinger residual test + muon cross-check (Candidate~3)}
\label{sec:eq-schwinger-residual}

Source:
\texttt{Mathematica\_Notebooks/Quantum\_Mechanics/r\_e\_schwinger\_residual\_test.wl}
(225 lines, scaffold for issue~\#66; Candidate-3 arc).

\paragraph{Test.}
Does the residual between the triangulated $r_e/r_0 = 0.4994205099128317$
(PR~\#62) and the closed-form Schwinger 1-loop value
$(2 - \alpha/(2\pi))/(4 + \alpha/\pi)$ track the higher-order QED corrections
to $g_e$?

\paragraph{Structural caveat (load-bearing).}
The triangulated $r_e$ is defined to make $g_r(r) = g_e^{\rm measured}$ to 16
sig figs (Pass~B fit). Therefore
$\Delta g_e^{\rm obs} \equiv g_e^{\rm measured} - g_e^{\rm Schwinger}$ is
\emph{exactly equal by construction} to the all-orders QED contributions
beyond Schwinger one-loop. The agreement measured in \S4 below is a
\emph{consistency check} that the triangulated value lies on the all-orders
QED curve in $r$-space, not an independent test.

\paragraph{Numerics (\S1--4).}
With $\alpha = 1/137.035999084$ (CODATA 2018):
\begin{center}
\begin{tabular}{ll}
\toprule
Quantity                                            & Value \\
\midrule
$\Delta g_e^{\rm obs} = g_e^{\rm meas} - g_e^{\rm Schwinger}$ & $\approx 3.515\times 10^{-6}$ \\
$\Delta g_e^{\rm pred}$ (KK + LR + KF, 2--4 loop)   & matches $\Delta g_e^{\rm obs}$ to $\sim 5\times 10^{-12}$ \\
\bottomrule
\end{tabular}
\end{center}
\wolframcheck{KK $C_2 = 0.328478965579193$, LR $C_3 = 1.181241456587$, KF
$C_4 \approx 1.9106$; signed sums per
$\Delta g_e^{(n)} = -2 \cdot a_e^{(n)}$.}

\paragraph{Tentative outcome-matrix branch: B-conditional.}
\begin{itemize}[topsep=2pt,itemsep=0pt]
  \item \textbf{B-confirmed} iff Tepper indicates the DRQM~I cutoff derivation
        produces the closed-form $(2 - \alpha/(2\pi))/(4 + \alpha/\pi)$ as an
        identity (with higher-order QED coming from framework loop
        corrections not yet derived in the campaign).
  \item \textbf{A-confirmed} iff the framework's cutoff derivation produces
        directly the triangulated value $0.4994205099128317$ (requires the
        cutoff to encode all-orders QED at the derivation level --- suspicious
        without a derived loop structure).
  \item \textbf{D-confirmed} iff neither, in which case the closed-form
        Schwinger match is contingent.
\end{itemize}

\paragraph{Muon cross-check (\S8) --- breaking the back-fit degeneracy.}
Applying the same $g_r$ formula to the muon: with FNAL Muon~$g-2$ 2023 final
$a_\mu^{\rm exp} = 1.16592059\times 10^{-3}$
($\sigma_{a_\mu} = 2.2\times 10^{-10}$) and
$a_e^{\rm exp} = 1.15965218073\times 10^{-3}$,
\[
\frac{r_e}{r_0^e} - \frac{r_\mu}{r_0^\mu} \;\approx\; 3.13\times 10^{-6}
\quad\equiv\quad \approx 57{,}000\;\text{muon-}\sigma
\]
\wolframcheck{Module evaluated 2026-05-26.}

\paragraph{Universal-closed-form-cutoff hypothesis: ruled out at $\sim 57{,}000\,\sigma$.}
The per-particle back-fit cutoffs differ by $3\times 10^{-6}$; equivalently
$a_\mu - a_e = 6.27\times 10^{-6}$ differs from zero at $\sim 28{,}500$
muon-$\sigma$. \textbf{However, the published DRQM~I does not make a
universal-cutoff prediction} --- the paper writes (III.23) for the muon and
proton with $r_\mu$, $r_p$ as separate parameters without numerical values.
The cross-check rules out any future \emph{closed-form universal} reading
(e.g.~from issue~\#54) that would predict particle-mass-independent cutoffs.

\paragraph{Implication for Outcome~C.}
Strengthens C-as-published: the framework's cutoff is necessarily per-particle
and empirically fit to each particle's measured anomaly; the closed-form
match for the electron is the back-fit cutoff that reproduces $a_e^{\rm exp}$;
the same logic applied to the muon would yield
$r_\mu/r_0^\mu = (2 - \alpha/(2\pi))/(4 + \alpha/\pi) + \delta_\mu$ with
$\delta_\mu$ from $a_\mu^{\rm exp} - \alpha/(2\pi) \sim 4.5\times 10^{-6}$,
similar size to the electron's deviation.

\paragraph{Crocco compliance.}
Numerical evaluation, Wolfram-MCP cell execution, and constant tabulation are
pragmatic AI. The structural caveat, the interpretation of the residual as
not-independent, and the formulation of the Tepper-discriminating question
are substantive AI.

\authorreview{(1)~Does DRQM~I~\S III.D derive a closed-form $r_e/r_0$ in
$\alpha$, or does it not? (2)~If yes, is the derived form the closed-form
$(2 - \alpha/(2\pi))/(4 + \alpha/\pi)$ (B-branch) or the triangulated value
directly (A-branch)? (3)~Is the muon cross-check correctly framed as a
constraint on \emph{future} closed-form derivations rather than a
falsification of the framework as written?}
